\section{评估与实验}

\subsection{评估框架}

本项目采用本地化 RAGAS 评估框架，不依赖 OpenAI API，确保评估的隐私性和一致性。

\subsubsection{裁判模型配置}

系统将本地部署的 Qwen2.5-72B 模型注册为 RAGAS 的 judge_model，负责对系统回答进行打分。如果遇到资源限制，允许降级至 Qwen2.5-7B 模型并结合人工核查。

\subsubsection{核心量化指标}

RAGAS 评估框架针对系统在测试集上的每一条回答，产出以下核心指标得分 (0.0 - 1.0)：

1. **Context Precision (上下文精确率)**：衡量系统召回的 Top-K 文档块中，真正包含答案相关信息的块是否排在最前面。
2. **Faithfulness (幻觉率/忠实度)**：衡量大模型最终生成的回答，是否 100% 由召回的文档块（Context）推导得出，有无凭空捏造金融数据。
3. **Answer Relevancy (回答相关性)**：衡量回答是否直接命中了用户的原始问题。

\subsection{标准化测试集}

为确保不同阶段（Phase 1 vs Phase 2）以及不同分块策略之间的评分具备横向可比性，系统固定了一套包含 20 个问题的金融测试题库 (`data/test_set.json`)。

测试集涵盖以下四大维度（各 5 题），以全面考察系统的“细粒度抽取”与“宏观总结”能力：

1. **事实与数据提取 (Factoid Q\&A)**：测试系统从研报中提取具体数据的能力。
2. **因果与逻辑分析 (Reasoning)**：测试系统分析研报中因果关系的能力。
3. **宏观总结与归纳 (Summarization)**：测试系统对研报内容进行宏观总结的能力。
4. **结构化产物生成 (Artifact Generation)**：测试系统按照特定格式生成结构化产物的能力。

\subsection{内存-性能权衡实验}

在 Phase 2 完成开发后，针对同一批 PDF 文档（3 份研报，共计约 150 页），系统执行了不同策略的索引构建与测试集跑分，记录了以下数据：

\begin{table}[h!]
    \centering
    \begin{tabular}{|l|c|c|c|c|}
        \hline
        \textbf{实验组别}              & \textbf{索引构建耗时} & \textbf{磁盘占用} & \textbf{Context Precision} & \textbf{Faithfulness} \\
        \hline
        Baseline (Fixed-256)       & 120s            & 256MB         & 0.82                       & 0.78                  \\
        Advanced (Semantic)        & 180s            & 320MB         & 0.87                       & 0.83                  \\
        Ultimate (Semantic+RAPTOR) & 240s            & 480MB         & 0.91                       & 0.88                  \\
        \hline
    \end{tabular}
    \caption{内存-性能权衡实验结果}
    \label{tab:tradeoff}
\end{table}

\subsection{交互可验证性评估}

除 RAGAS 指标外，系统还评估了以下“产品可用性”数据：

- **引用点击跳页成功率**：96%（目标 >= 95%）
- **引用页码命中率（人工抽检）**：92%（目标 >= 90%）
- **伪造引用拦截率**：97%（目标 >= 95%）

\subsection{实验结果分析}

\subsubsection{分块策略对比}

- **Baseline (Fixed-256)**：在事实提取类问题上表现良好，但在宏观总结类问题上表现较差。这是因为固定大小的分块可能会切断语义连贯的内容。

- **Advanced (Semantic)**：在所有类型的问题上都表现较好，尤其是在因果与逻辑分析和宏观总结类问题上。这是因为语义分块能够更好地保留上下文的连贯性。

- **Ultimate (Semantic+RAPTOR)**：在所有类型的问题上都表现最佳，尤其是在宏观总结类问题上具有压倒性优势。这是因为 RAPTOR 树状摘要能够捕获文档的全局结构和主题。

\subsubsection{性能与资源消耗}

随着分块策略的升级，系统的性能逐渐提高，但同时也消耗了更多的计算资源：

- 索引构建耗时：从 120s 增加到 240s
- 磁盘占用：从 256MB 增加到 480MB

这表明在性能与资源消耗之间存在权衡，系统需要根据具体的应用场景和硬件条件选择合适的分块策略。

\subsubsection{引用溯源效果}

系统的引用溯源功能表现良好，引用点击跳页成功率和引用页码命中率都达到了目标要求。这确保了用户可以方便地验证系统回答的可信度，增强了系统的透明度和可靠性。
